\section{Introduction}

\subsection{Motivation and Problem Statement}

Transformer architectures are the dominant paradigm for sequence modeling in natural language processing \citep{vaswani_attention_2017}, computer vision, and computational biology, with models such as BERT \citep{devlin_bert_2019} and GPT establishing the standard for representation learning. The field has since diversified rapidly: recurrent and retentive blocks, sparse attention, gated mechanisms, latent projections, memory-augmented layers, fast-weight rules, and continuous-depth formulations each address a different limitation of dense attention, while the optimizer landscape has expanded beyond AdamW with variance-reduction, memory-efficient, schedule-free, second-order, and low-rank methods \citep{child_sparse_transformer_2019,sun_retentive_2023,gu_mamba_2023,behrouz_titans_2025,ramapuram_theory_2024,xie_adan_2022,pagliardini_ademamix_2024,yuan_mars_2024,shazeer_adafactor_2018,zhao_galore_2024,defazio_road_2024,gupta_shampoo_2018,vyas_soap_2024}.

The practical consequence is a fragmented ecosystem. Comparing sequence mixers or optimizers under identical training conditions requires re-implementing each variant, wiring it into a training pipeline, and maintaining the result across framework changes. This effort hinders reproducibility, slows scientific progress, and raises the barrier to entry for new researchers.

\subsection{Proposal and Contributions}

We propose \textbf{Frankenstein Transformer} (\url{https://github.com/erickfmm/frankenstein-transformer}), a configuration-driven toolkit that treats architecture, optimizer, and training choices as a single validated document. The project provides:

\begin{enumerate}[leftmargin=*]
    \item \textbf{Schema-first design}: a strict YAML contract supporting forty-three sequence mixers across eight families and twenty-three optimizer families, enforced before training begins.
    \item \textbf{Architecture-agnostic training}: one pipeline for dense baselines, recurrent and retentive blocks, sparse patterns, gated mechanisms, latent compression, conditional memory, fast-weight attention, and geometric-field attention, with encoder and decoder modes \citep{zhu2026mixtureofdepthsattention,cheng2026conditionalmemoryscalablelookup,xie_mhc_2025,zhang_falcon_2026,pisoni_ssog_2026}.
    \item \textbf{Optimizer routing}: prefixed hyperparameter groups that control embeddings, normalization layers, attention blocks, and remaining parameters independently for each optimizer family \citep{zhu2025apollosgdlikememoryadamwlevel}.
    \item \textbf{End-to-end workflows}: training (including sentence-embedding objectives), deployment and export (quantized, standard, HuggingFace Transformers, and BitNet GGUF formats), inference, and a schema-driven web builder inspired by SBERT-style workflows \citep{reimers_sentence-bert_2019}.
    \item \textbf{Vision support}: a Vision Transformer model class covering masked patch prediction, classification, and segmentation under the same mixer and optimizer infrastructure.
    \item \textbf{Validation tooling}: unit tests, YAML preset validation, and continuous integration across supported Python versions.
\end{enumerate}

The command surface is \texttt{frankenstein-transformer}, with four subcommands: \texttt{train}, \texttt{deploy}, \texttt{infer}, and \texttt{web-server}. The last launches a Streamlit configuration builder that renders schema metadata as forms, validates input live, previews the resulting YAML, and generates the equivalent CLI command, serving both as a configuration aid and as documentation of the available design space.

\subsection{Reading Guide}

This document is organized as a technical reference:
\begin{enumerate}[leftmargin=*]
    \item \textbf{Prerequisites}: Appendix~\ref{sec:annex-tutorial} introduces transformers and attention for readers new to the field.
    \item \textbf{Configuration contract}: Section~\ref{sec:config-centric-architecture} describes the schema and Section~\ref{sec:schema-validation} its validation rules.
    \item \textbf{Architecture selection}: Section~\ref{sec:attention-families} compares sequence-mixer families; Appendices~\ref{sec:annex-transformer}, \ref{sec:annex-latent}, \ref{sec:annex-sparse}, and \ref{sec:annex-gated} collect the supporting literature, with normalization, BitNet, and conditional-memory details in Appendix~\ref{sec:annex-norm-bitnet-mod} and manifold-constrained residuals in Appendix~\ref{sec:annex-mhc}.
    \item \textbf{Optimization}: Section~\ref{sec:optimizer-families} summarizes optimizer routing; Appendix~\ref{sec:annex-optimizers} provides the full family analysis.
    \item \textbf{Deployment and inference}: Sections~\ref{sec:quantization} and~\ref{sec:sbert} cover quantized deployment and sentence embeddings.
\end{enumerate}
